% ===== Block B: PDF p.4–9 | 原文 §2.1（Definition of the moduli spaces） | 唯一编号公式 2.1–2.16 =====
% [NOTE: 本块自 §2.1 标题之后开始；§2 与 §2.1 两级标题由 Lead 写，不在本块内。]
% [原文 p.4]

\noindent\textbf{术语.} 我们研究的主要对象是光滑、连通、紧的复一维流形，简称为曲线或黎曼曲面，它们带有 $n$ 个标记的不同点（见 M. Bertola 的讲义 [Ber]）。这些对象将记作
\begin{equation}
  (\Sigma,p_1,\dots,p_n)\,.
  \tag{2.1}
\end{equation}
% [Tutor Note: (2.1) 的记号把曲线与其 $n$ 个标记点当作一个整体对象 $(\Sigma,p_1,\dots,p_n)$；亏格 $g$ 是底层实曲面的不变量，而 $n$ 个点是有标号且互异的，这两条正是后文 $\cM_{g,n}$ 与稳定曲线定义中反复出现的数据。核对 $n=0$ 的约定：此时记号退化为单个 $\Sigma$，$2g-2+n=2g-2$。]

每条紧复曲线都带有底层的实二维可定向、无边的紧曲面的结构，而后者由它的亏格 $g$ 唯一刻画。

我们主要的例子将是球面（亏格 $0$）与环面（亏格 $1$）。球面作为黎曼曲面在同构意义下只有唯一的结构，即复射影直线 $\PP^1$。亏格为 $0$ 的复曲线称为\emph{有理曲线}。$\PP^1$ 的自同构群是
\begin{equation}
  \mathrm{PSL}(2,\C)=\left\{\begin{pmatrix}a&b\\c&d\end{pmatrix}\ \middle|\
    \begin{array}{c}[a:b:c:d]\in\PP^3\\ ad-bc\ne 0\end{array}\right\}
  \tag{2.2}
\end{equation}
它作用为
\begin{equation}
  \begin{pmatrix}a&b\\c&d\end{pmatrix}\!\cdot z=\frac{az+b}{cz+d}\,.
  \tag{2.3}
\end{equation}
% [Tutor Note: (2.2) 中的条件 $[a:b:c:d]\in\PP^3$ 不是多余的：$ad-bc$ 是矩阵元素的二次齐次式，故 $ad-bc\ne 0$ 在整体缩放 $(a,b,c,d)\mapsto\lambda(a,b,c,d)$ 下不变，可以在 $\PP^3$ 中陈述；而集合 (2.2) 就是 $\mathrm{GL}(2,\C)/\C^*\cong\mathrm{PSL}(2,\C)$，即 $\PP^3$ 中由该条件定义的开集。核对：$a=d=1$、$b=c=0$ 给出恒等变换，与 (2.3) 中 $z\mapsto z$ 一致。]
% [Tutor Note: (2.3) 是矩阵左乘给出的 Möbius 变换。用三个特殊值核对矩阵与分式的对应：$z=0\mapsto b/d$，$z=\infty\mapsto a/c$，$z=1\mapsto(a+b)/(c+d)$；它们在 $(a,b,c,d)$ 整体乘 $\lambda$ 时都不变，这也与 (2.2) 用 $\PP^3$ 参数化一致。]
% [Tutor Note: 由单值化定理，亏格 $0$ 的紧黎曼曲面必共形同构于 $\PP^1=\C\cup\{\infty\}$，这正是正文“至多同构唯一”的依据。也正因为这种唯一性，$\cM_{0,0}$ 只有一个点；而该点上的自同构群是无限的 $\mathrm{PSL}(2,\C)$，所以它“没有好的几何结构”。]

至于亏格 $1$，环面上的每个黎曼曲面结构在同构意义下都可作为商 $\C/\Lambda$ 得到。这里 $\Lambda$ 是一个格，即形如
\begin{equation}
  \Lambda=\{n_1\omega_1+n_2\omega_2\mid n_1,n_2\in\Z\}
  \tag{2.4}
\end{equation}
的加法群，其中 $\omega_1,\omega_2\in\C$ 在实数上线性无关。亏格为 $1$ 的复曲线称为\emph{椭圆曲线}。

正如引言中所讨论的，我们关心的是固定亏格 $g$、带 $n$ 个标记点的黎曼曲面的模空间（特别是，我们希望使“在这样的空间上作积分”有意义：即二维量子引力的路径积分）。

\figomit{2}{实二维可定向紧曲面的若干例子。}

% [原文 p.5]
\begin{definition}[2.1]
  \emph{模空间} $\cM_{g,n}$ 是亏格 $g$、带 $n$ 个标记点的黎曼曲面的同构类构成的集合：
  \begin{equation}
    \cM_{g,n}=\left\{\begin{array}{c}
      \text{亏格 }g\text{、带 }n\text{ 个标记点的黎曼曲面}
    \end{array}\right\}\Big/\,\text{同构}
    \tag{2.5}
  \end{equation}
  两个对象 $(\Sigma,p_1,\dots,p_n)$ 与 $(\Sigma',p'_1,\dots,p'_n)$ 之间的同构，指的是一个保持标记点的双全纯映射 $\phi:\Sigma\to\Sigma'$，即 $\phi(p_i)=p'_i$。
\end{definition}
% [Tutor Note: 定义 2.1 只把 $\cM_{g,n}$ 定义为集合，右下角的 “iso.” 表示按同构取商；它带有复轨形结构要等到 $2g-2+n>0$（定理 2.2）。核对极端情形：$g=0,n=3$ 时 (2.5) 是单点集，与 (2.6) 第一行一致。]

上述定义完全没问题，但我们希望赋予它更多的结构。回忆一下，我们的目标是在黎曼曲面的模空间上讨论积分，因此具有类似流形的结构将是件好事。事实证明，这里有很多几何，但并不像流形的几何那样好。主要原因是黎曼曲面有自同构。最简单的例子是 $\PP^1$，它的自同构群是无限群 $\mathrm{PSL}(2,\C)$。由于在积分中我们想模去对称群，一个无限的自同构群是个坏消息。换句话说，$\cM_{0,0}$ 并没有好的几何结构。不过，有一种摆脱自同构的办法：标记（至少三个）点。

\begin{exercise}[2.1]
  \begin{enumerate}[label=(\arabic*)]
    \item 考虑带三个标记点 $(\PP^1,p_1,p_2,p_3)$ 的亏格 $0$ 曲线。求出把 $(\PP^1,p_1,p_2,p_3)$ 映到 $(\PP^1,0,1,\infty)$ 的（唯一的）$g\in\mathrm{PSL}(2,\C)$。
    \item 考虑带四个标记点 $(\PP^1,p_1,p_2,p_3,p_4)$ 的亏格 $0$ 曲线。第 (1) 问中求出的群元素 $g\in\mathrm{PSL}(2,\C)$ 把 $(\PP^1,p_1,p_2,p_3,p_4)$ 映到 $(\PP^1,0,1,\infty,t)$。把 $t$ 表示成 $p_1,p_2,p_3,p_4$ 的函数\footnote{这就是所谓的\emph{交比}（cross-ratio），它的定义可上溯到远古（或许欧几里得时已有），Pappus 曾研究过它，并指出了它关键的不变性。}。
  \end{enumerate}
\end{exercise}
% [Tutor Note: 练习 2.1(2) 的 $t$ 就是交比。取 $t=\frac{(p_4-p_1)(p_2-p_3)}{(p_4-p_3)(p_2-p_1)}$（在 $\PP^1=\C\cup\{\infty\}$ 中按极限理解），则 $p_4$ 分别等于 $p_1,p_2,p_3$ 时 $t=0,1,\infty$，故四点互异等价于 $t\notin\{0,1,\infty\}$，这正是 (2.6) 第二行。注意答案依赖于归一化的次序：换一种把三个点送到 $0,1,\infty$ 的方式，$t$ 会被一个分式线性变换替换（仍满足 $t\notin\{0,1,\infty\}$）。]

上面的练习表明
\begin{align}
  \cM_{0,3}&=\{(\PP^1,0,1,\infty)\}=\{*\}\,,\nonumber\\
  \cM_{0,4}&=\{(\PP^1,0,1,\infty,t)\mid t\ne 0,1,\infty\}=\PP^1\setminus\{0,1,\infty\}\,.
  \tag{2.6}
\end{align}

把上面的分析推广，可以证明当 $n\ge 3$ 时
\begin{equation}
  \cM_{0,n}=\left\{(t_1,\dots,t_{n-3})\in(\PP^1\setminus\{0,1,\infty\})^{n-3}\ \middle|\ t_i\ne t_j\right\}.
  \tag{2.7}
\end{equation}
这就赋予 $\cM_{0,n}$ 以好的几何结构。
% [Tutor Note: (2.7) 的由来：由练习 2.1(1)，三个互异标记点可被 $\mathrm{PSL}(2,\C)$ 唯一地送到 $0,1,\infty$；剩下 $n-3$ 个点记为 $t_1,\dots,t_{n-3}$，它们必须落在 $\PP^1\setminus\{0,1,\infty\}$ 内且两两不同，反过来这些数据唯一确定同构类。核对 $n=3$：空元组只有一个，退化为 (2.6) 的 $\{*\}$。]

另一个自同构群无限的反例是椭圆曲线 $E$：$\Aut(E)$ 含有与 $E$ 本身同构的子群，它通过平移作用。同样地，我们可以通过标记一个点来摆脱自同构（这里指平移）。若 $E=\C/\Lambda$，标记点的一个自然选择是 $\Lambda\subset\C$ 的像，即环面上的单位元。于是 $\cM_{1,1}=\{\mbox{格}\}/\C^*$，其中 $\C^*$ 通过重新缩放作用。为理解这个商，让我们固定 $\Lambda$ 的一组基 $(\omega_1,\omega_2)$。把 $\Lambda$ 乘以 $1/\omega_1$，我们得到一组等价格的基 $(1,\tau)$，其中 $\tau$ 属于上半平面 $\HH$。选取同一个格的另一组基，即用格的基变换群 $\mathrm{SL}(2,\Z)$ 作用，我们得到另一个点 $\tau'\in\HH$。于是我们发现
\begin{equation}
  \cM_{1,1}=\HH/\mathrm{SL}(2,\Z)\,.
  \tag{2.8}
\end{equation}

% [原文 p.6]
\figomit{3}{模空间 $\cM_{1,1}$。弧 $AB$ 与 $AB'$ 以及半直线 $BC$ 与 $B'C'$ 分别被等同。}

该商的一个基本域如图 3 所示。在把弧 $AB\sim AB'$ 与半直线 $BC\sim B'C'$ 等同之后，我们看到 $\cM_{1,1}$ 拓扑上就是 $\PP^1\setminus\{\infty\}$。然而格具有非平凡的自同构。事实上，矩阵 $-\mathrm{Id}$ 在 $\HH$ 上作用平凡，因此 $\cM_{1,1}$ 上每个点的自同构群至少含有 $\Z_2$ 作为子群。这个对合称为带标记椭圆曲线的\emph{超椭圆对合}。若把椭圆曲线写成形如
\begin{equation}
  E:\quad y^2=x^3+ax+b\,,
  \tag{2.9}
\end{equation}
的三次多项式方程（的紧化），那么超椭圆对合就只是映射 $y\mapsto -y$。
% [Tutor Note: (2.8) 写成 $\mathrm{SL}(2,\Z)$ 而模去的是 $\mathrm{PSL}(2,\Z)$，二者一致：$-I\in\mathrm{SL}(2,\Z)$ 在 $\HH$ 上作用平凡，这正对应后文“矩阵 $-\mathrm{Id}$ 在 $\HH$ 上作用平凡”。核对生成元：$S:\tau\mapsto -1/\tau$ 与 $T:\tau\mapsto\tau+1$ 生成 $\mathrm{SL}(2,\Z)$，基本域即图 3 中经典区域 $\{|\tau|\ge 1,\ |\operatorname{Re}\tau|\le 1/2\}$。]
% [Tutor Note: (2.9) 用的是 Weierstrass 形式；当判别式（至多差非零常数因子）$4a^3+27b^2\ne 0$ 时它给出光滑椭圆曲线。后文 (2.15) 的 Legendre 形式 $y^2=x(x-1)(x-t)$ 是同一个对象的另一种写法，$t\to 0,1$ 时两根碰撞、曲线退化。]
% [Tutor Note: (2.9) 里的超椭圆对合 $y\mapsto -y$ 与 $\PP^1$ 上的 $x$ 坐标相容：$(x,y)\mapsto x$ 是 $\PP^1$ 的二重覆盖，分歧点恰为 $x^3+ax+b$ 的三个根以及 $x=\infty$（共 4 个，与 $g=1$ 的 Riemann–Hurwitz 计数一致）；$\Z_2$ 交换两个叶，只在分歧点处不动。]

由 A. Hurwitz 的一个定理可知，任何满足 $2g-2+n>0$ 的黎曼曲面，其自同构群都有限。这样的数对 $(g,n)$ 称为\emph{稳定的}。反过来，每条满足 $2g-2+n\le 0$ 的黎曼曲面都有无限多个保持标记点的自同构。换言之：
\begin{equation}
  \Aut(\Sigma_g,p_1,\dots,p_n)\ \mbox{有限}
  \Longleftrightarrow -\chi=2g-2+n>0\,.
  \tag{2.10}
\end{equation}
这就排除了把 $\cM_{0,0},\cM_{0,1},\cM_{0,2}$ 以及 $\cM_{1,0}$ 定义为好的几何空间的可能。（它们虽然仍可看作集合，但这样做的用处有限。）
% [Tutor Note: (2.10) 中的 $\chi$ 指 $\chi(\Sigma_g)=2-2g$，故 $-\chi=2g-2+n$ 是带 $n$ 个标记点的穿孔曲面的负 Euler 示性数。核对稳定性条件：$(g,n)=(0,3)$ 给 $1>0$、$(1,1)$ 给 $1>0$（稳定），$(0,2)$ 给 $0$、$(1,0)$ 给 $0$（不稳定，正是正文列出的坏例子）。]

事实上，完全可以刻画基本域中每个点 $\tau$ 的自同构群（见图 4）：
\begin{itemize}
  \item 对 $\tau=\ee^{\pi\ii/3}=\frac{1+\ii\sqrt3}{2}$（对应六角格），自同构群是 $\Z_6$；
  \item 对 $\tau=\ee^{\pi\ii/2}=\ii$（对应正方格），自同构群是 $\Z_4$；
  \item 对基本域中任何其他 $\tau$，自同构群是 $\Z_2$。
\end{itemize}
% [Tutor Note: 核对数值：$\ee^{\pi\ii/3}=\cos 60^\circ+\ii\sin 60^\circ=\frac{1+\ii\sqrt3}{2}$，$\ee^{\pi\ii/2}=\ii$。图 4 第三个示例标的是 $\tau=\ee^{5\pi\ii/24}$，它并不在经典基本域内（$\operatorname{Re}\ee^{5\pi\ii/24}\approx 0.79>1/2$），图中只是“一般格”的示例，与正文“基本域中其他 $\tau$”并不矛盾。]

% [原文 p.7]
\figomit{4}{格的自同构群。}

从现在起，我们总假设 $2g-2+n>0$。在这种情形下情况是好的，但还没有好到极致：仍存在具有非平凡自同构群的曲线，正如 $\cM_{1,1}$ 的例子所示。尽管如此，自同构群的有限性使我们能够把黎曼曲面的模空间视为一个轨形（orbifold）。

\begin{theorem}[2.2]
  对 $2g-2+n>0$，模空间 $\cM_{g,n}$ 是连通的、光滑的复轨形，其维数为
  \begin{equation}
    \dim(\cM_{g,n})=3g-3+n\,.
    \tag{2.11}
  \end{equation}
\end{theorem}
% [Tutor Note: (2.11) 的三个核对：$g=0,n=3$ 给 $0$（单点），$g=0,n=4$ 给 $1$（与 (2.6) 的 $\PP^1\setminus\{0,1,\infty\}$ 一致），$g=1,n=1$ 给 $1$（与 (2.8) 的 $\HH/\mathrm{SL}(2,\Z)$ 一致）。一般地这对应“$3g-3$ 个复结构模参数加上 $n$ 个标记点位置”。]

光滑复轨形的定义颇有些技术性，但它与光滑复流形的定义十分相似。关键区别在于：局部地，轨形看起来像 $\C^d/G$ 的开子集，其中 $G$ 是一个有限群，全纯地作用。一个值得记住的简单例子是整体商 $\C/\Z_m$，其中 $\Z_m$ 通过角度 $2\pi/m$ 的旋转作用。在上面这个定理中，“光滑”是在这种轨形意义下说的：模空间 $\cM_{g,n}$ 局部地以光滑复流形对有限群作用的商为模型。用实际的话说，这意味着 $\cM_{g,n}$ 在相差有限个商奇点的意义下是光滑流形。

\begin{exercise}[2.2]
  对熟悉 Riemann–Roch 与 Riemann–Hurwitz 的读者，请说服自己：$\cM_g=\cM_{g,0}$ 的复维数是 $3g-3$。这一结果 Riemann 本人已经知道，他还创造了“模空间”（moduli space）一词（源自拉丁文 modus，意为度量）：

  % [Tutor Note: 引文中的德文是 Riemann 的原话（底本未注明出处，按内容应出自其 1857 年关于 Abelian 函数的论文），此处按其直译。术语对照：Moduln 即“模数”（不是“模空间”），$p$ 即现在的亏格 $g$，“$2p+1$ fach zusammenhängend” 即“$2p+1$ 重连通”（沿 $2p$ 条闭曲线切开后化为单连通），Verzweigungswerthe 即分支值。]
  \begin{quote}\small
    然而，上述关于一类 $2p+1$ 重连通的代数函数之模数个数的确定，只有在下述前提下才成立：存在 $2\mu-p+1$ 个分支值，它们是函数 $\xi$ 中任意常数的彼此独立的函数。这一前提仅在 $p>1$ 时成立，而此时模数的个数才等于 $3p-3$；当 $p=1$ 时则等于 $1$。不过，对这些模数的直接研究之所以变得困难，原因在于任意常数出现在 $\xi$ 中的方式。因此，为了确定模数的个数，人们在由具有相同分支的 $2p+1$ 重连通代数函数所构成的系统中，取作独立变动的量并不是其中某个函数，而是这样一个函数的一个处处有限的积分。
  \end{quote}

  % [原文 p.8]
  为此，考虑二元组 $(\Sigma,f)$ 的模空间，其中 $\Sigma$ 是亏格 $g$ 的黎曼曲面，$f$ 是从 $\Sigma$ 到 $\PP^1$ 的 $d$ 次全纯映射（即 $\Sigma$ 上的亚纯函数）。这样的空间有时称为 Hurwitz 空间，记作 $\cH_{g,d}$。请用两种不同的方式计算它的维数。
  \begin{itemize}
    \item $\cH_{g,d}$ 的维数等于 $\cM_g$ 的维数，即黎曼曲面 $\Sigma$ 的“形变参数个数”，加上函数 $f$ 的“形变参数个数”。后者用 Riemann–Roch 计算。
    \item 用 Riemann–Hurwitz 直接计算 $\cH_{g,d}$ 的维数。
  \end{itemize}
  由此得出 $\dim\cM_g=3g-3$。

  \textbf{提示.} 考虑 $d\gg 1$。
\end{exercise}
% [Tutor Note: 练习 2.2 的两种算法必须给出同一维数：用 Riemann–Roch 算 $f$ 的形变参数时需 $d\gg 1$，此时 $h^0(\Sigma,f^*\cO_{\PP^1}(1))=d+1-g$（$h^1$ 项消失）；而由 Riemann–Hurwitz，$d$ 次映射的分支点数为 $2d+2g-2$。两相比较即得 $\dim\cM_g=3g-3$。另外，原文 “Hint.” 前印有一个小灯泡图标，此处以“提示.”代之，未另加图标。]

这一局部描述的一个重要推论是：许多从流形中熟悉的微分几何概念，例如微分形式与积分，都自然地推广到轨形上。特别地，在复轨形上谈论积分是有意义的。以 $\C/\Z_m$ 为例，给定一个在角度 $2\pi/m$ 的旋转下不变的函数 $f:\C\to\C$，我们可以定义
\begin{equation}
  \int_{\C/\Z_m}f(z,\bar z)\,\dd z\,\dd\bar z
  =\frac{1}{|\Z_m|}\int_{\C}f(z,\bar z)\,\dd z\,\dd\bar z\,.
  \tag{2.12}
\end{equation}
% [Tutor Note: (2.12) 就是轨形积分的定义：先在覆盖 $\C$ 上积分，再除以 $\Z_m$ 的阶。核对 $f\equiv 1$：左边是 $\C/\Z_m$ 的“轨形面积”，右边是 $\C$ 中任一基本域的面积，二者一致；这也说明为何要求 $f$ 在 $2\pi/m$ 的旋转下不变，否则右边会依赖基本域的选取。]

对流形成立的大多数结果都（在作适当修正后）推广到轨形上。下面是 Euler 示性数的一个例子。

\begin{exercise}[2.3]
  轨形 $X$ 的 Euler 示性数定义为
  \begin{equation}
    \chi(X)=\sum_{G}\frac{\chi(X_G)}{|G|}\,,
    \tag{2.13}
  \end{equation}
  其中 $X_G$ 是自同构群为 $G$ 的点构成的子集。证明 $\chi(\cM_{1,1})=-\frac{1}{12}$。该公式推广为著名的 Harer–Zagier 公式 [HZ86]：
  \begin{equation}
    \chi(\cM_{g,n})=(1-2g)_{n-1}\,\zeta(1-2g)\,,
    \tag{2.14}
  \end{equation}
  其中 $(x)_m$ 表示 Pochhammer 符号（或下降阶乘），$\zeta(x)$ 是 Riemann zeta 函数。有趣的是，Harer 与 Zagier 最初的计算使用的是矩阵模型方法。
\end{exercise}
% [Tutor Note: 用 (2.14) 核对前面的断言：$g=1,n=1$ 时 $(1-2g)_0=1$、$\zeta(-1)=-\frac{1}{12}$，恰得 $\chi(\cM_{1,1})=-\frac{1}{12}$。注意 $g=0$ 时 $\zeta(1-2g)=\zeta(1)$ 有极点，故该式只适用于 $g\ge 1$；这里 $(x)_m=x(x-1)\cdots(x-m+1)$ 是下降阶乘，$(x)_0=1$。]

虽然轨形上的积分是良定义的，但还有另一个潜在的问题需要处理：非紧性。非紧性问题从 $\cM_{0,4}$ 或 $\cM_{1,1}$ 的例子中已经可以看出。后者在拓扑上是 $\PP^1\setminus\{\infty\}$，其中缺失的无穷远点正是非紧性的来源。通过考虑椭圆曲线族
\begin{equation}
  E_t:\quad y^2=x(x-1)(x-t)\,,\qquad t\in(0,1)\,,
  \tag{2.15}
\end{equation}
我们可以看出这个极限点是如何几何地实现的。在 $t\to 0$ 或 $1$ 的极限下，黎曼曲面 $E_t$ 变成退化的。例如，当 $t\to 0$ 时我们得到 $y^2=x^2(x-1)$，它在 $x=0$ 附近看起来像两条复直线 $y=\pm x$ 的并。这意味着在 $x=0$ 处有两个相汇的分支，也就是所谓的结点奇点，而曲面 $E_0$ 将如图 5 所示。换句话说，$\cM_{1,1}$ 的极限点不再是一个环面，而是一个捏缩环面。
% [Tutor Note: 核对“局部像两条复直线之并”：在 $x=0$ 附近有 $y=\pm x\sqrt{x-1}=\pm\ii x\sqrt{1-x}$，取局部全纯坐标 $X=\ii x$、$Y=y/\sqrt{1-x}$，方程化为 $Y^2=X^2$，即两条光滑分支 $Y=\pm X$ 在原点横截相汇——这正是结点的局部模型。故正文的 “$y=\pm x$” 应理解为在适当的局部全纯坐标下。]

% [原文 p.9]
\figomit{5}{一个捏缩环面。}

为了使得在非紧空间上作积分有意义，我们有两种可能。第一种是只考虑在极限点处具有某种衰减的函数或微分形式。第二种选择是把所关心的空间适当地紧化，然后只考虑这种紧化上的正则函数或正则微分形式。我们将采用第二条路线。事实证明，对于 $\cM_{g,n}$，只需添入带结点的黎曼曲面，就足以得到好的紧化。

\begin{definition}[2.3]
  亏格 $g$、带 $n$ 个有标号标记点 $p_1,\dots,p_n$ 的\emph{稳定黎曼曲面}，是一个可能奇异的、紧的、连通的复曲线 $\Sigma$，使得：
  \begin{itemize}
    \item 把所有结点光滑化后所得的曲面的亏格为 $g$（见图 6）；
    \item $\Sigma$ 的奇点只有结点；
    \item 标记点互不相同，且不与结点重合；并且
    \item $(\Sigma,p_1,\dots,p_n)$ 只有有限个自同构。
  \end{itemize}
\end{definition}
% [Tutor Note: 定义 2.3 的最后一条与“每个不可约分支上 $2g(v)-2+n(v)>0$”等价（见本块末段）：$n(v)$ 把结点在该分支上的原像也计入标记点，于是每个分支的自同构群有限，整体自同构群才有限。核对不稳定例：两条 $\PP^1$ 以一个结点相连且无标记点，则每个分支上 $2\cdot 0-2+2=0$，存在 $\C^*$ 型的连续自同构族。注意原文用 “smoothening”（非标准拼法，通常作 smoothing），本译统一作“光滑化”。]

于是我们可以定义一个参数化稳定黎曼曲面的同构类的模空间，它常被称为 Deligne–Mumford 模空间 [DM69]：
\begin{equation}
  \Mc_{g,n}=\left\{\begin{array}{c}
    \text{亏格 }g\text{、带 }n\text{ 个标记点的稳定黎曼曲面}
  \end{array}\right\}\Big/\,\text{同构}
  \tag{2.16}
\end{equation}
% [Tutor Note: 与 (2.5) 相比，(2.16) 只是把“黎曼曲面”换成“稳定黎曼曲面”，因此 $\cM_{g,n}\subset\Mc_{g,n}$。本块只给出集合层面的定义；$\Mc_{g,n}$ 的紧性以及 $\cM_{g,n}$ 在其中的开稠密性，属于下一块的定理 2.4。]

\figomit{6}{一个奇异黎曼曲面的光滑化与正规化。由光滑化读出 $(g,n)$；由正规化读出稳定性条件。}

% [原文 p.10]（本节末段自 p.9 跨页至此，按“开始者写完”归本块）
上述定义中的最后一个条件可以重新表述如下。设 $\Sigma_1,\dots,\Sigma_k$ 为把每个结点处的分支分开后所得曲面的连通分支（这一过程称为正规化，见图 6）。设 $g(v)$ 是 $\Sigma_v$ 的亏格，$n(v)$ 是特殊点的个数，即标记点与各结点在 $\Sigma_i$ 上的原像。于是，“有限自同构”条件成立当且仅当对所有的 $v$ 都有 $2g(v)-2+n(v)>0$。
% [Tutor Note: 末句的 $\Sigma_i$ 与同句前面的 $\Sigma_v$ 不一致：按上下文，指标 $v$ 遍历连通分支 $\Sigma_1,\dots,\Sigma_k$，故此处应为 $\Sigma_v$；这一处原文笔误已在 ERRATA 中登记，正文照原样排印。]
